% Supplementary material for main.tex: the complete feature inventory and the
% per-condition tree splits. Both appendices are GENERATED by make_appendix.py.
% Build main.tex first (with --keep-intermediates) so that xr can resolve the
% cross-references into the main paper from main.aux.
\documentclass[lettersize,journal]{IEEEtran}

\usepackage{booktabs}
\usepackage{amsmath,amssymb}
\usepackage{float}
\usepackage{array}
\usepackage{xr}
\externaldocument{main}
\usepackage[colorlinks=true,linkcolor=black,citecolor=black,urlcolor=blue]{hyperref}
\usepackage{siunitx}
\sisetup{detect-weight=true,detect-family=true}
\DeclareSIUnit{\point}{point}
\DeclareSIUnit{\sample}{sample}

\newcommand{\ttf}[1]{\texttt{\footnotesize #1}}

\begin{document}

\title{Supplementary Material for\\
Audio Codec Effects on Feature Attribution\\
in Speech Emotion Recognition}
\author{Akshath~R}
\markboth{IEEE Transactions on Affective Computing --- Supplementary Material}%
{R: Audio Codec Effects on Feature Attribution in SER}
\maketitle

Section, equation and figure numbers without an S prefix refer to the main
paper.

% GENERATED BY make_appendix.py -- DO NOT EDIT BY HAND
\appendices

\section{Complete Feature Inventory}\label{app:features}

Every stimulus is reduced to 436 named columns. Frame-level contours are summarised by eight order statistics (mean, standard deviation, minimum, maximum, median, interquartile range, skewness, excess kurtosis). Descriptors marked with a shorter list are summarised only by those statistics, and scalars are already clip-level. Naming is \ttf{<family>\_<descriptor>\_<statistic>} throughout, which is what makes an attribution over a column a claim about an acoustic quantity.

Analysis frames are \SI{512}{point} FFT, \SI{400}{sample} window, \SI{160}{sample} hop at \SI{16}{\kilo\hertz}, which gives \SI{25}{\milli\second} windows every \SI{10}{\milli\second}. Praat measures use a \SIrange{75}{500}{\hertz} pitch range.

\subsection{Descriptors and their definitions}

\begin{table}[H]
\caption{Spectral descriptors. 114 columns of 436. ``all 8'' denotes the eight order statistics of Section~\ref{sec:features}: mean, standard deviation, minimum, maximum, median, IQR, skewness and excess kurtosis.}
\label{tab:app-spectral}
\centering\scriptsize\setlength{\tabcolsep}{3pt}
\begin{tabular}{@{}>{\raggedright\arraybackslash}p{2.05cm}>{\raggedright\arraybackslash}p{3.85cm}>{\raggedright\arraybackslash}p{1.20cm}r@{}}
\toprule
Descriptor & Definition & Stats & $n$\\
\midrule
Spectral centroid & First moment of the power spectrum. It is the frequency about which spectral energy is balanced. Correlates with perceived brightness. & all 8 & 8\\
Spectral bandwidth & Power-weighted second moment about the centroid. It measures how widely energy is spread in frequency. & all 8 & 8\\
Spectral roll-off (85\,\%) & Frequency below which 85\,\% of the frame's spectral energy lies. & all 8 & 8\\
Spectral roll-off (95\,\%) & As above at 95\,\%. Directly sensitive to codec low-pass behaviour. & all 8 & 8\\
Spectral flatness & Geometric mean of the magnitude spectrum divided by its arithmetic mean. It approaches 1 for noise-like frames and 0 for tonal ones, and it collapses when an encoder sets masked bins to zero. & all 8 & 8\\
Zero-crossing rate & Fraction of adjacent sample pairs with a sign change. It is a cheap proxy for voicing and frication. & all 8 & 8\\
Frame RMS energy & Root-mean-square magnitude per frame, which gives the loudness contour. & all 8 & 8\\
Spectral flux & Euclidean norm of the frame-to-frame difference of the magnitude spectrum. It measures how fast the spectrum is changing. & all 8 & 8\\
Spectral entropy & Shannon entropy of the power spectrum normalised to its own bin count, $-\sum_k P_k\ln P_k/\ln K$ with $P_k$ the share of frame power in bin $k$ of $K$. It is a scale-free measure of how evenly energy is distributed across frequency. & all 8 & 8\\
Spectral slope & Least-squares slope of magnitude regressed on frequency, per frame. This is the spectral tilt. & all 8 & 8\\
Band-energy ratios (8 bands) & Fraction of frame energy in each of 0--500, 500--1000, 1000--2000, 2000--3000, 3000--4000, 4000--5000, 5000--6000, 6000--8000\,Hz. Energy-normalised, so these describe spectral \emph{shape} rather than level. & mean, std & 16\\
High-frequency survival ratios & Fraction of frame energy above 4 and above 6\,kHz. The most direct probes of lossy-codec low-pass behaviour in the inventory. & mean, std & 4\\
Spectral contrast (7 sub-bands) & Per-band difference between spectral peaks and valleys in decibels. Band 6 covers 6--8\,kHz and is the single most codec-displaced column in the table (Section~\ref{sec:signal}). & mean, std & 14\\
\bottomrule
\end{tabular}
\end{table}

\begin{table}[H]
\caption{Cepstral descriptors. 240 columns of 436. ``all 8'' denotes the eight order statistics of Section~\ref{sec:features}: mean, standard deviation, minimum, maximum, median, IQR, skewness and excess kurtosis.}
\label{tab:app-mfcc}
\centering\scriptsize\setlength{\tabcolsep}{3pt}
\begin{tabular}{@{}>{\raggedright\arraybackslash}p{2.05cm}>{\raggedright\arraybackslash}p{3.85cm}>{\raggedright\arraybackslash}p{1.20cm}r@{}}
\toprule
Descriptor & Definition & Stats & $n$\\
\midrule
$\Delta$MFCC (20) & First-order temporal derivative of each MFCC, over a 9-frame window. & mean, std & 40\\
$\Delta\Delta$MFCC (20) & Second-order temporal derivative of each MFCC. & mean, std & 40\\
MFCC (20) & Discrete cosine transform of the log mel-filterbank energies, which is the standard SER front end. MFCC-0 tracks overall log energy, so its derivatives measure how fast loudness changes. & all 8 & 160\\
\bottomrule
\end{tabular}
\end{table}

\begin{table}[H]
\caption{Prosodic and voice-quality descriptors (Praat). 56 columns of 436. ``all 8'' denotes the eight order statistics of Section~\ref{sec:features}: mean, standard deviation, minimum, maximum, median, IQR, skewness and excess kurtosis.}
\label{tab:app-prosody}
\centering\scriptsize\setlength{\tabcolsep}{3pt}
\begin{tabular}{@{}>{\raggedright\arraybackslash}p{2.05cm}>{\raggedright\arraybackslash}p{3.85cm}>{\raggedright\arraybackslash}p{1.20cm}r@{}}
\toprule
Descriptor & Definition & Stats & $n$\\
\midrule
$F_0$ slope (mean $|\cdot|$) & Mean absolute first difference of the voiced $F_0$ contour, per second. It is the rate of pitch movement. & scalar & 1\\
Fundamental frequency $F_0$ & Praat autocorrelation pitch over voiced frames only. & all 8 & 8\\
Jitter (4 variants) & Cycle-to-cycle variation in glottal period: local, RAP, PPQ5, DDP. A perturbation measure of the voice source. & scalar & 4\\
Shimmer (6 variants) & Cycle-to-cycle variation in glottal amplitude: local, local\,dB, APQ3, APQ5, APQ11, DDA. APQ11 needs eleven consecutive periods and is the one column in the table with non-trivial missingness (11.3\,\%). & scalar & 6\\
Harmonics-to-noise ratio & Ratio of periodic to aperiodic energy in decibels. It correlates with hoarseness and breathiness. & mean, std, min, max, median & 5\\
Formant bandwidths $B_1$--$B_3$ & Median $-3$\,dB bandwidth of the first three formants (Burg estimator). & median & 1\\
Formants $F_1$--$F_5$ & Burg-estimator formant frequencies, which describe the vocal-tract filter. $F_2$ movement is the strongest single correlate of disgust in this corpus. & mean, std, median & 17\\
Intensity & Praat short-term intensity contour in decibels. Its standard deviation is among the top-ranked features for every model family in the study. & all 8 & 8\\
Voiced fraction & Proportion of frames with a defined $F_0$. & scalar & 1\\
Pause fraction & Complement of the voiced fraction. By definition it equals $1 -$ voiced fraction ($r = -1.00$). & scalar & 1\\
Voiced-segment count & Number of maximal runs of voiced frames. & scalar & 1\\
Voiced-segment rate & Voiced segments per second, which is a proxy for speech rate. & scalar & 1\\
Mean voiced-segment duration & Mean length in seconds of a voiced run. & scalar & 1\\
CPPS & Smoothed cepstral peak prominence, which is the standard acoustic correlate of breathy phonation. Drives both sadness and neutrality in every attribution method applied here. & scalar & 1\\
\bottomrule
\end{tabular}
\end{table}

\begin{table}[H]
\caption{Chroma descriptors. 24 columns of 436. ``all 8'' denotes the eight order statistics of Section~\ref{sec:features}: mean, standard deviation, minimum, maximum, median, IQR, skewness and excess kurtosis.}
\label{tab:app-chroma}
\centering\scriptsize\setlength{\tabcolsep}{3pt}
\begin{tabular}{@{}>{\raggedright\arraybackslash}p{2.05cm}>{\raggedright\arraybackslash}p{3.85cm}>{\raggedright\arraybackslash}p{1.20cm}r@{}}
\toprule
Descriptor & Definition & Stats & $n$\\
\midrule
Chroma (12 pitch classes) & Energy folded onto the twelve semitone pitch classes. Tuning is pinned to 0 rather than estimated, because an estimated tuning could itself shift under compression and become part of the effect being measured. & mean, std & 24\\
\bottomrule
\end{tabular}
\end{table}

\begin{table}[H]
\caption{Global descriptors. 2 columns of 436. ``all 8'' denotes the eight order statistics of Section~\ref{sec:features}: mean, standard deviation, minimum, maximum, median, IQR, skewness and excess kurtosis.}
\label{tab:app-global}
\centering\scriptsize\setlength{\tabcolsep}{3pt}
\begin{tabular}{@{}>{\raggedright\arraybackslash}p{2.05cm}>{\raggedright\arraybackslash}p{3.85cm}>{\raggedright\arraybackslash}p{1.20cm}r@{}}
\toprule
Descriptor & Definition & Stats & $n$\\
\midrule
Clip duration & Decoded duration in seconds. Rises by 32.68\,ms under AAC encoder priming, uniformly across every clip. & scalar & 1\\
Peak amplitude & Maximum absolute sample value after loudness normalisation. & scalar & 1\\
\bottomrule
\end{tabular}
\end{table}

\subsection{Enumerated column names}

The complete inventory, in the order the extractor emits it.

\begin{table*}[!tp]
\caption{Spectral descriptors: all 114 column names.}
\label{tab:app-list-spectral}
\centering\tiny\setlength{\tabcolsep}{3pt}
\begin{tabular}{@{}llll@{}}
\toprule
\texttt{spectral\_centroid\_mean} & \texttt{spectral\_rolloff95\_iqr} & \texttt{spectral\_flux\_min} & \texttt{spectral\_band\_2000\_3000\_std}\\
\texttt{spectral\_centroid\_std} & \texttt{spectral\_rolloff95\_skew} & \texttt{spectral\_flux\_max} & \texttt{spectral\_band\_3000\_4000\_mean}\\
\texttt{spectral\_centroid\_min} & \texttt{spectral\_rolloff95\_kurt} & \texttt{spectral\_flux\_median} & \texttt{spectral\_band\_3000\_4000\_std}\\
\texttt{spectral\_centroid\_max} & \texttt{spectral\_flatness\_mean} & \texttt{spectral\_flux\_iqr} & \texttt{spectral\_band\_4000\_5000\_mean}\\
\texttt{spectral\_centroid\_median} & \texttt{spectral\_flatness\_std} & \texttt{spectral\_flux\_skew} & \texttt{spectral\_band\_4000\_5000\_std}\\
\texttt{spectral\_centroid\_iqr} & \texttt{spectral\_flatness\_min} & \texttt{spectral\_flux\_kurt} & \texttt{spectral\_band\_5000\_6000\_mean}\\
\texttt{spectral\_centroid\_skew} & \texttt{spectral\_flatness\_max} & \texttt{spectral\_entropy\_mean} & \texttt{spectral\_band\_5000\_6000\_std}\\
\texttt{spectral\_centroid\_kurt} & \texttt{spectral\_flatness\_median} & \texttt{spectral\_entropy\_std} & \texttt{spectral\_band\_6000\_8000\_mean}\\
\texttt{spectral\_bandwidth\_mean} & \texttt{spectral\_flatness\_iqr} & \texttt{spectral\_entropy\_min} & \texttt{spectral\_band\_6000\_8000\_std}\\
\texttt{spectral\_bandwidth\_std} & \texttt{spectral\_flatness\_skew} & \texttt{spectral\_entropy\_max} & \texttt{spectral\_hf\_ratio\_4000\_mean}\\
\texttt{spectral\_bandwidth\_min} & \texttt{spectral\_flatness\_kurt} & \texttt{spectral\_entropy\_median} & \texttt{spectral\_hf\_ratio\_4000\_std}\\
\texttt{spectral\_bandwidth\_max} & \texttt{spectral\_zcr\_mean} & \texttt{spectral\_entropy\_iqr} & \texttt{spectral\_hf\_ratio\_6000\_mean}\\
\texttt{spectral\_bandwidth\_median} & \texttt{spectral\_zcr\_std} & \texttt{spectral\_entropy\_skew} & \texttt{spectral\_hf\_ratio\_6000\_std}\\
\texttt{spectral\_bandwidth\_iqr} & \texttt{spectral\_zcr\_min} & \texttt{spectral\_entropy\_kurt} & \texttt{spectral\_contrast\_0\_mean}\\
\texttt{spectral\_bandwidth\_skew} & \texttt{spectral\_zcr\_max} & \texttt{spectral\_slope\_mean} & \texttt{spectral\_contrast\_0\_std}\\
\texttt{spectral\_bandwidth\_kurt} & \texttt{spectral\_zcr\_median} & \texttt{spectral\_slope\_std} & \texttt{spectral\_contrast\_1\_mean}\\
\texttt{spectral\_rolloff85\_mean} & \texttt{spectral\_zcr\_iqr} & \texttt{spectral\_slope\_min} & \texttt{spectral\_contrast\_1\_std}\\
\texttt{spectral\_rolloff85\_std} & \texttt{spectral\_zcr\_skew} & \texttt{spectral\_slope\_max} & \texttt{spectral\_contrast\_2\_mean}\\
\texttt{spectral\_rolloff85\_min} & \texttt{spectral\_zcr\_kurt} & \texttt{spectral\_slope\_median} & \texttt{spectral\_contrast\_2\_std}\\
\texttt{spectral\_rolloff85\_max} & \texttt{spectral\_rms\_mean} & \texttt{spectral\_slope\_iqr} & \texttt{spectral\_contrast\_3\_mean}\\
\texttt{spectral\_rolloff85\_median} & \texttt{spectral\_rms\_std} & \texttt{spectral\_slope\_skew} & \texttt{spectral\_contrast\_3\_std}\\
\texttt{spectral\_rolloff85\_iqr} & \texttt{spectral\_rms\_min} & \texttt{spectral\_slope\_kurt} & \texttt{spectral\_contrast\_4\_mean}\\
\texttt{spectral\_rolloff85\_skew} & \texttt{spectral\_rms\_max} & \texttt{spectral\_band\_0\_500\_mean} & \texttt{spectral\_contrast\_4\_std}\\
\texttt{spectral\_rolloff85\_kurt} & \texttt{spectral\_rms\_median} & \texttt{spectral\_band\_0\_500\_std} & \texttt{spectral\_contrast\_5\_mean}\\
\texttt{spectral\_rolloff95\_mean} & \texttt{spectral\_rms\_iqr} & \texttt{spectral\_band\_500\_1000\_mean} & \texttt{spectral\_contrast\_5\_std}\\
\texttt{spectral\_rolloff95\_std} & \texttt{spectral\_rms\_skew} & \texttt{spectral\_band\_500\_1000\_std} & \texttt{spectral\_contrast\_6\_mean}\\
\texttt{spectral\_rolloff95\_min} & \texttt{spectral\_rms\_kurt} & \texttt{spectral\_band\_1000\_2000\_mean} & \texttt{spectral\_contrast\_6\_std}\\
\texttt{spectral\_rolloff95\_max} & \texttt{spectral\_flux\_mean} & \texttt{spectral\_band\_1000\_2000\_std} & \\
\texttt{spectral\_rolloff95\_median} & \texttt{spectral\_flux\_std} & \texttt{spectral\_band\_2000\_3000\_mean} & \\
\bottomrule
\end{tabular}
\end{table*}

\begin{table*}[!tp]
\caption{Cepstral descriptors: all 240 column names.}
\label{tab:app-list-mfcc}
\centering\tiny\setlength{\tabcolsep}{3pt}
\begin{tabular}{@{}llll@{}}
\toprule
\texttt{mfcc\_00\_mean} & \texttt{mfcc\_07\_median} & \texttt{mfcc\_15\_mean} & \texttt{mfcc\_d1\_05\_mean}\\
\texttt{mfcc\_00\_std} & \texttt{mfcc\_07\_iqr} & \texttt{mfcc\_15\_std} & \texttt{mfcc\_d1\_05\_std}\\
\texttt{mfcc\_00\_min} & \texttt{mfcc\_07\_skew} & \texttt{mfcc\_15\_min} & \texttt{mfcc\_d2\_05\_mean}\\
\texttt{mfcc\_00\_max} & \texttt{mfcc\_07\_kurt} & \texttt{mfcc\_15\_max} & \texttt{mfcc\_d2\_05\_std}\\
\texttt{mfcc\_00\_median} & \texttt{mfcc\_08\_mean} & \texttt{mfcc\_15\_median} & \texttt{mfcc\_d1\_06\_mean}\\
\texttt{mfcc\_00\_iqr} & \texttt{mfcc\_08\_std} & \texttt{mfcc\_15\_iqr} & \texttt{mfcc\_d1\_06\_std}\\
\texttt{mfcc\_00\_skew} & \texttt{mfcc\_08\_min} & \texttt{mfcc\_15\_skew} & \texttt{mfcc\_d2\_06\_mean}\\
\texttt{mfcc\_00\_kurt} & \texttt{mfcc\_08\_max} & \texttt{mfcc\_15\_kurt} & \texttt{mfcc\_d2\_06\_std}\\
\texttt{mfcc\_01\_mean} & \texttt{mfcc\_08\_median} & \texttt{mfcc\_16\_mean} & \texttt{mfcc\_d1\_07\_mean}\\
\texttt{mfcc\_01\_std} & \texttt{mfcc\_08\_iqr} & \texttt{mfcc\_16\_std} & \texttt{mfcc\_d1\_07\_std}\\
\texttt{mfcc\_01\_min} & \texttt{mfcc\_08\_skew} & \texttt{mfcc\_16\_min} & \texttt{mfcc\_d2\_07\_mean}\\
\texttt{mfcc\_01\_max} & \texttt{mfcc\_08\_kurt} & \texttt{mfcc\_16\_max} & \texttt{mfcc\_d2\_07\_std}\\
\texttt{mfcc\_01\_median} & \texttt{mfcc\_09\_mean} & \texttt{mfcc\_16\_median} & \texttt{mfcc\_d1\_08\_mean}\\
\texttt{mfcc\_01\_iqr} & \texttt{mfcc\_09\_std} & \texttt{mfcc\_16\_iqr} & \texttt{mfcc\_d1\_08\_std}\\
\texttt{mfcc\_01\_skew} & \texttt{mfcc\_09\_min} & \texttt{mfcc\_16\_skew} & \texttt{mfcc\_d2\_08\_mean}\\
\texttt{mfcc\_01\_kurt} & \texttt{mfcc\_09\_max} & \texttt{mfcc\_16\_kurt} & \texttt{mfcc\_d2\_08\_std}\\
\texttt{mfcc\_02\_mean} & \texttt{mfcc\_09\_median} & \texttt{mfcc\_17\_mean} & \texttt{mfcc\_d1\_09\_mean}\\
\texttt{mfcc\_02\_std} & \texttt{mfcc\_09\_iqr} & \texttt{mfcc\_17\_std} & \texttt{mfcc\_d1\_09\_std}\\
\texttt{mfcc\_02\_min} & \texttt{mfcc\_09\_skew} & \texttt{mfcc\_17\_min} & \texttt{mfcc\_d2\_09\_mean}\\
\texttt{mfcc\_02\_max} & \texttt{mfcc\_09\_kurt} & \texttt{mfcc\_17\_max} & \texttt{mfcc\_d2\_09\_std}\\
\texttt{mfcc\_02\_median} & \texttt{mfcc\_10\_mean} & \texttt{mfcc\_17\_median} & \texttt{mfcc\_d1\_10\_mean}\\
\texttt{mfcc\_02\_iqr} & \texttt{mfcc\_10\_std} & \texttt{mfcc\_17\_iqr} & \texttt{mfcc\_d1\_10\_std}\\
\texttt{mfcc\_02\_skew} & \texttt{mfcc\_10\_min} & \texttt{mfcc\_17\_skew} & \texttt{mfcc\_d2\_10\_mean}\\
\texttt{mfcc\_02\_kurt} & \texttt{mfcc\_10\_max} & \texttt{mfcc\_17\_kurt} & \texttt{mfcc\_d2\_10\_std}\\
\texttt{mfcc\_03\_mean} & \texttt{mfcc\_10\_median} & \texttt{mfcc\_18\_mean} & \texttt{mfcc\_d1\_11\_mean}\\
\texttt{mfcc\_03\_std} & \texttt{mfcc\_10\_iqr} & \texttt{mfcc\_18\_std} & \texttt{mfcc\_d1\_11\_std}\\
\texttt{mfcc\_03\_min} & \texttt{mfcc\_10\_skew} & \texttt{mfcc\_18\_min} & \texttt{mfcc\_d2\_11\_mean}\\
\texttt{mfcc\_03\_max} & \texttt{mfcc\_10\_kurt} & \texttt{mfcc\_18\_max} & \texttt{mfcc\_d2\_11\_std}\\
\texttt{mfcc\_03\_median} & \texttt{mfcc\_11\_mean} & \texttt{mfcc\_18\_median} & \texttt{mfcc\_d1\_12\_mean}\\
\texttt{mfcc\_03\_iqr} & \texttt{mfcc\_11\_std} & \texttt{mfcc\_18\_iqr} & \texttt{mfcc\_d1\_12\_std}\\
\texttt{mfcc\_03\_skew} & \texttt{mfcc\_11\_min} & \texttt{mfcc\_18\_skew} & \texttt{mfcc\_d2\_12\_mean}\\
\texttt{mfcc\_03\_kurt} & \texttt{mfcc\_11\_max} & \texttt{mfcc\_18\_kurt} & \texttt{mfcc\_d2\_12\_std}\\
\texttt{mfcc\_04\_mean} & \texttt{mfcc\_11\_median} & \texttt{mfcc\_19\_mean} & \texttt{mfcc\_d1\_13\_mean}\\
\texttt{mfcc\_04\_std} & \texttt{mfcc\_11\_iqr} & \texttt{mfcc\_19\_std} & \texttt{mfcc\_d1\_13\_std}\\
\texttt{mfcc\_04\_min} & \texttt{mfcc\_11\_skew} & \texttt{mfcc\_19\_min} & \texttt{mfcc\_d2\_13\_mean}\\
\texttt{mfcc\_04\_max} & \texttt{mfcc\_11\_kurt} & \texttt{mfcc\_19\_max} & \texttt{mfcc\_d2\_13\_std}\\
\texttt{mfcc\_04\_median} & \texttt{mfcc\_12\_mean} & \texttt{mfcc\_19\_median} & \texttt{mfcc\_d1\_14\_mean}\\
\texttt{mfcc\_04\_iqr} & \texttt{mfcc\_12\_std} & \texttt{mfcc\_19\_iqr} & \texttt{mfcc\_d1\_14\_std}\\
\texttt{mfcc\_04\_skew} & \texttt{mfcc\_12\_min} & \texttt{mfcc\_19\_skew} & \texttt{mfcc\_d2\_14\_mean}\\
\texttt{mfcc\_04\_kurt} & \texttt{mfcc\_12\_max} & \texttt{mfcc\_19\_kurt} & \texttt{mfcc\_d2\_14\_std}\\
\texttt{mfcc\_05\_mean} & \texttt{mfcc\_12\_median} & \texttt{mfcc\_d1\_00\_mean} & \texttt{mfcc\_d1\_15\_mean}\\
\texttt{mfcc\_05\_std} & \texttt{mfcc\_12\_iqr} & \texttt{mfcc\_d1\_00\_std} & \texttt{mfcc\_d1\_15\_std}\\
\texttt{mfcc\_05\_min} & \texttt{mfcc\_12\_skew} & \texttt{mfcc\_d2\_00\_mean} & \texttt{mfcc\_d2\_15\_mean}\\
\texttt{mfcc\_05\_max} & \texttt{mfcc\_12\_kurt} & \texttt{mfcc\_d2\_00\_std} & \texttt{mfcc\_d2\_15\_std}\\
\texttt{mfcc\_05\_median} & \texttt{mfcc\_13\_mean} & \texttt{mfcc\_d1\_01\_mean} & \texttt{mfcc\_d1\_16\_mean}\\
\texttt{mfcc\_05\_iqr} & \texttt{mfcc\_13\_std} & \texttt{mfcc\_d1\_01\_std} & \texttt{mfcc\_d1\_16\_std}\\
\texttt{mfcc\_05\_skew} & \texttt{mfcc\_13\_min} & \texttt{mfcc\_d2\_01\_mean} & \texttt{mfcc\_d2\_16\_mean}\\
\texttt{mfcc\_05\_kurt} & \texttt{mfcc\_13\_max} & \texttt{mfcc\_d2\_01\_std} & \texttt{mfcc\_d2\_16\_std}\\
\texttt{mfcc\_06\_mean} & \texttt{mfcc\_13\_median} & \texttt{mfcc\_d1\_02\_mean} & \texttt{mfcc\_d1\_17\_mean}\\
\texttt{mfcc\_06\_std} & \texttt{mfcc\_13\_iqr} & \texttt{mfcc\_d1\_02\_std} & \texttt{mfcc\_d1\_17\_std}\\
\texttt{mfcc\_06\_min} & \texttt{mfcc\_13\_skew} & \texttt{mfcc\_d2\_02\_mean} & \texttt{mfcc\_d2\_17\_mean}\\
\texttt{mfcc\_06\_max} & \texttt{mfcc\_13\_kurt} & \texttt{mfcc\_d2\_02\_std} & \texttt{mfcc\_d2\_17\_std}\\
\texttt{mfcc\_06\_median} & \texttt{mfcc\_14\_mean} & \texttt{mfcc\_d1\_03\_mean} & \texttt{mfcc\_d1\_18\_mean}\\
\texttt{mfcc\_06\_iqr} & \texttt{mfcc\_14\_std} & \texttt{mfcc\_d1\_03\_std} & \texttt{mfcc\_d1\_18\_std}\\
\texttt{mfcc\_06\_skew} & \texttt{mfcc\_14\_min} & \texttt{mfcc\_d2\_03\_mean} & \texttt{mfcc\_d2\_18\_mean}\\
\texttt{mfcc\_06\_kurt} & \texttt{mfcc\_14\_max} & \texttt{mfcc\_d2\_03\_std} & \texttt{mfcc\_d2\_18\_std}\\
\texttt{mfcc\_07\_mean} & \texttt{mfcc\_14\_median} & \texttt{mfcc\_d1\_04\_mean} & \texttt{mfcc\_d1\_19\_mean}\\
\texttt{mfcc\_07\_std} & \texttt{mfcc\_14\_iqr} & \texttt{mfcc\_d1\_04\_std} & \texttt{mfcc\_d1\_19\_std}\\
\texttt{mfcc\_07\_min} & \texttt{mfcc\_14\_skew} & \texttt{mfcc\_d2\_04\_mean} & \texttt{mfcc\_d2\_19\_mean}\\
\texttt{mfcc\_07\_max} & \texttt{mfcc\_14\_kurt} & \texttt{mfcc\_d2\_04\_std} & \texttt{mfcc\_d2\_19\_std}\\
\bottomrule
\end{tabular}
\end{table*}

\begin{table*}[!tp]
\caption{Prosodic and voice-quality descriptors (Praat): all 56 column names.}
\label{tab:app-list-prosody}
\centering\tiny\setlength{\tabcolsep}{3pt}
\begin{tabular}{@{}llll@{}}
\toprule
\texttt{prosody\_f0\_mean} & \texttt{prosody\_jitter\_local} & \texttt{prosody\_hnr\_median} & \texttt{prosody\_f4\_std}\\
\texttt{prosody\_f0\_std} & \texttt{prosody\_jitter\_rap} & \texttt{prosody\_f1\_mean} & \texttt{prosody\_f4\_median}\\
\texttt{prosody\_f0\_min} & \texttt{prosody\_jitter\_ppq5} & \texttt{prosody\_f1\_std} & \texttt{prosody\_f5\_mean}\\
\texttt{prosody\_f0\_max} & \texttt{prosody\_jitter\_ddp} & \texttt{prosody\_f1\_median} & \texttt{prosody\_f5\_std}\\
\texttt{prosody\_f0\_median} & \texttt{prosody\_shimmer\_local} & \texttt{prosody\_f1\_bw\_median} & \texttt{prosody\_f5\_median}\\
\texttt{prosody\_f0\_iqr} & \texttt{prosody\_shimmer\_localdb} & \texttt{prosody\_f2\_mean} & \texttt{prosody\_intensity\_mean}\\
\texttt{prosody\_f0\_skew} & \texttt{prosody\_shimmer\_apq3} & \texttt{prosody\_f2\_std} & \texttt{prosody\_intensity\_std}\\
\texttt{prosody\_f0\_kurt} & \texttt{prosody\_shimmer\_apq5} & \texttt{prosody\_f2\_median} & \texttt{prosody\_intensity\_min}\\
\texttt{prosody\_voiced\_fraction} & \texttt{prosody\_shimmer\_apq11} & \texttt{prosody\_f2\_bw\_median} & \texttt{prosody\_intensity\_max}\\
\texttt{prosody\_f0\_slope\_abs\_mean} & \texttt{prosody\_shimmer\_dda} & \texttt{prosody\_f3\_mean} & \texttt{prosody\_intensity\_median}\\
\texttt{prosody\_n\_voiced\_segments} & \texttt{prosody\_hnr\_mean} & \texttt{prosody\_f3\_std} & \texttt{prosody\_intensity\_iqr}\\
\texttt{prosody\_voiced\_rate\_per\_s} & \texttt{prosody\_hnr\_std} & \texttt{prosody\_f3\_median} & \texttt{prosody\_intensity\_skew}\\
\texttt{prosody\_mean\_voiced\_seg\_s} & \texttt{prosody\_hnr\_min} & \texttt{prosody\_f3\_bw\_median} & \texttt{prosody\_intensity\_kurt}\\
\texttt{prosody\_pause\_fraction} & \texttt{prosody\_hnr\_max} & \texttt{prosody\_f4\_mean} & \texttt{prosody\_cpps}\\
\bottomrule
\end{tabular}
\end{table*}

\begin{table}[!tp]
\caption{Chroma descriptors: all 24 column names.}
\label{tab:app-list-chroma}
\centering\tiny\setlength{\tabcolsep}{3pt}
\begin{tabular}{@{}lll@{}}
\toprule
\texttt{chroma\_00\_mean} & \texttt{chroma\_04\_mean} & \texttt{chroma\_08\_mean}\\
\texttt{chroma\_00\_std} & \texttt{chroma\_04\_std} & \texttt{chroma\_08\_std}\\
\texttt{chroma\_01\_mean} & \texttt{chroma\_05\_mean} & \texttt{chroma\_09\_mean}\\
\texttt{chroma\_01\_std} & \texttt{chroma\_05\_std} & \texttt{chroma\_09\_std}\\
\texttt{chroma\_02\_mean} & \texttt{chroma\_06\_mean} & \texttt{chroma\_10\_mean}\\
\texttt{chroma\_02\_std} & \texttt{chroma\_06\_std} & \texttt{chroma\_10\_std}\\
\texttt{chroma\_03\_mean} & \texttt{chroma\_07\_mean} & \texttt{chroma\_11\_mean}\\
\texttt{chroma\_03\_std} & \texttt{chroma\_07\_std} & \texttt{chroma\_11\_std}\\
\bottomrule
\end{tabular}
\end{table}

\begin{table}[!tp]
\caption{Global descriptors: all 2 column names.}
\label{tab:app-list-global}
\centering\tiny\setlength{\tabcolsep}{3pt}
\begin{tabular}{@{}lll@{}}
\toprule
\texttt{duration\_s} & \texttt{peak\_amplitude} & \\
\bottomrule
\end{tabular}
\end{table}

\clearpage


% GENERATED BY make_appendix.py -- DO NOT EDIT BY HAND
\section{Per-Condition Tree Splits}\label{app:splits}

The tables below list every internal node down to depth 4 of the entropy tree of Section~\ref{sec:trees}. The tree was fitted independently on each condition, using the same actors, the same items and the same hyper-parameters. The \emph{Path} column reads from the root, and \ttf{L} denotes the branch taken when the parent's test is satisfied. Entropies are given in bits. The value $n$ is the raw sample count, and the information gain is computed on the class-weighted counts (Eq.~\ref{eq:ig}).

A path missing from a table is one the tree stopped splitting, because a child would have fallen below the 25-sample leaf minimum. Comparing the same path across the four tables is the node-level comparison summarised in Fig.~\ref{fig:depth} of the main paper.

\begin{table}[H]
\caption{Entropy tree fitted on \textbf{Uncompressed WAV}, all internal nodes to depth 4.}
\label{tab:app-splits-ref}
\centering\scriptsize\setlength{\tabcolsep}{1.5pt}
\begin{tabular}{@{}llrrrr@{}}
\toprule
Path & Feature & Threshold & $H$ & IG & $n$\\
\midrule
\textit{root} & \texttt{prosody\_intensity\_std} & 8.2184 & 2.5850 & 0.3090 & 4905\\
\texttt{L} & \texttt{mfcc\_d1\_00\_std} & 3.9965 & 2.3849 & 0.1415 & 3065\\
\texttt{R} & \texttt{mfcc\_d1\_00\_std} & 6.6630 & 2.0891 & 0.1797 & 1840\\
\texttt{LL} & \texttt{spectral\_band\_0\_500\_mean} & 0.8626 & 1.9382 & 0.0685 & 1059\\
\texttt{LR} & \texttt{prosody\_f0\_mean} & 220.5886 & 2.3995 & 0.0990 & 2006\\
\texttt{RL} & \texttt{prosody\_f0\_median} & 161.8673 & 2.2633 & 0.1141 & 847\\
\texttt{RR} & \texttt{prosody\_intensity\_min} & 42.7554 & 1.6055 & 0.0925 & 993\\
\texttt{LLL} & \texttt{prosody\_intensity\_median} & 64.3901 & 2.1110 & 0.0763 & 503\\
\texttt{LLR} & \texttt{prosody\_f0\_mean} & 174.4501 & 1.6481 & 0.0727 & 556\\
\texttt{LRL} & \texttt{spectral\_centroid\_iqr} & 209.5395 & 2.3328 & 0.0682 & 1781\\
\texttt{LRR} & \texttt{prosody\_f0\_max} & 372.4767 & 2.0340 & 0.1797 & 225\\
\texttt{RLL} & \texttt{spectral\_centroid\_iqr} & 310.5503 & 2.3335 & 0.1579 & 254\\
\texttt{RLR} & \texttt{prosody\_f0\_min} & 149.0330 & 2.0683 & 0.0778 & 593\\
\texttt{RRL} & \texttt{mfcc\_03\_max} & 46.1135 & 1.2265 & 0.0747 & 542\\
\texttt{RRR} & \texttt{spectral\_contrast\_2\_mean} & 12.1160 & 1.8550 & 0.0915 & 451\\
\texttt{LLLL} & \texttt{mfcc\_05\_iqr} & 8.7058 & 2.3273 & 0.1863 & 139\\
\texttt{LLLR} & \texttt{mfcc\_07\_mean} & 1.0004 & 1.9225 & 0.0759 & 364\\
\texttt{LLRL} & \texttt{mfcc\_00\_mean} & -220.3358 & 1.6344 & 0.0676 & 391\\
\texttt{LLRR} & \texttt{prosody\_intensity\_iqr} & 5.8890 & 1.4344 & 0.1343 & 165\\
\texttt{LRLL} & \texttt{mfcc\_d1\_03\_std} & 1.3160 & 2.2210 & 0.0810 & 564\\
\texttt{LRLR} & \texttt{duration\_s} & 2.6526 & 2.2848 & 0.0690 & 1217\\
\texttt{LRRL} & \texttt{mfcc\_d1\_00\_std} & 5.5015 & 2.2359 & 0.2615 & 119\\
\texttt{LRRR} & \texttt{mfcc\_d2\_05\_std} & 0.6661 & 1.4206 & 0.2061 & 106\\
\texttt{RLLL} & \texttt{spectral\_contrast\_1\_std} & 3.1735 & 2.2781 & 0.1619 & 146\\
\texttt{RLLR} & \texttt{duration\_s} & 2.4524 & 2.0339 & 0.1987 & 108\\
\texttt{RLRL} & \texttt{mfcc\_03\_std} & 12.0396 & 2.1028 & 0.0993 & 385\\
\texttt{RLRR} & \texttt{prosody\_intensity\_std} & 10.4187 & 1.7818 & 0.1437 & 208\\
\texttt{RRLL} & \texttt{spectral\_entropy\_mean} & 0.4682 & 1.5107 & 0.1241 & 305\\
\texttt{RRLR} & \texttt{mfcc\_d1\_00\_std} & 8.0503 & 0.6898 & 0.0951 & 237\\
\texttt{RRRL} & \texttt{mfcc\_d2\_00\_std} & 3.2530 & 2.1332 & 0.2152 & 95\\
\texttt{RRRR} & \texttt{mfcc\_03\_max} & 40.8027 & 1.6633 & 0.1040 & 356\\
\bottomrule
\end{tabular}
\end{table}

\begin{table}[H]
\caption{Entropy tree fitted on \textbf{MP3 @ 64 kbit/s}, all internal nodes to depth 4.}
\label{tab:app-splits-mp3-64}
\centering\scriptsize\setlength{\tabcolsep}{1.5pt}
\begin{tabular}{@{}llrrrr@{}}
\toprule
Path & Feature & Threshold & $H$ & IG & $n$\\
\midrule
\textit{root} & \texttt{prosody\_intensity\_std} & 8.2210 & 2.5850 & 0.3090 & 4905\\
\texttt{L} & \texttt{mfcc\_d1\_00\_std} & 4.0175 & 2.3849 & 0.1398 & 3065\\
\texttt{R} & \texttt{mfcc\_d1\_00\_std} & 6.6883 & 2.0891 & 0.1800 & 1840\\
\texttt{LL} & \texttt{spectral\_band\_0\_500\_mean} & 0.8627 & 1.9524 & 0.0669 & 1080\\
\texttt{LR} & \texttt{prosody\_f0\_mean} & 221.5304 & 2.3997 & 0.1002 & 1985\\
\texttt{RL} & \texttt{prosody\_f0\_median} & 161.8687 & 2.2615 & 0.1129 & 873\\
\texttt{RR} & \texttt{prosody\_intensity\_min} & 42.3291 & 1.5890 & 0.0914 & 967\\
\texttt{LLL} & \texttt{prosody\_intensity\_median} & 63.9287 & 2.1257 & 0.0752 & 515\\
\texttt{LLR} & \texttt{prosody\_f0\_mean} & 175.0993 & 1.6634 & 0.0762 & 565\\
\texttt{LRL} & \texttt{spectral\_band\_0\_500\_std} & 0.2524 & 2.3332 & 0.0688 & 1771\\
\texttt{LRR} & \texttt{mfcc\_02\_std} & 14.7187 & 2.0072 & 0.1668 & 214\\
\texttt{RLL} & \texttt{spectral\_centroid\_iqr} & 310.5299 & 2.3306 & 0.1592 & 257\\
\texttt{RLR} & \texttt{spectral\_band\_0\_500\_mean} & 0.7316 & 2.0707 & 0.0767 & 616\\
\texttt{RRL} & \texttt{spectral\_band\_3000\_4000\_mean} & 0.0146 & 1.2123 & 0.0774 & 532\\
\texttt{RRR} & \texttt{spectral\_contrast\_2\_mean} & 12.2081 & 1.8440 & 0.1046 & 435\\
\texttt{LLLL} & \texttt{mfcc\_05\_iqr} & 8.7206 & 2.3261 & 0.1391 & 148\\
\texttt{LLLR} & \texttt{spectral\_flatness\_mean} & 0.0041 & 1.9388 & 0.0744 & 367\\
\texttt{LLRL} & \texttt{mfcc\_00\_mean} & -223.6980 & 1.6513 & 0.0753 & 399\\
\texttt{LLRR} & \texttt{prosody\_intensity\_iqr} & 5.8541 & 1.4319 & 0.1347 & 166\\
\texttt{LRLL} & \texttt{prosody\_hnr\_median} & 3.7622 & 2.2810 & 0.0703 & 935\\
\texttt{LRLR} & \texttt{mfcc\_d2\_00\_std} & 2.4823 & 2.2461 & 0.0777 & 836\\
\texttt{LRRL} & \texttt{mfcc\_06\_mean} & -1.1106 & 1.4472 & 0.3107 & 65\\
\texttt{LRRR} & \texttt{mfcc\_16\_std} & 4.4168 & 2.0110 & 0.1969 & 149\\
\texttt{RLLL} & \texttt{mfcc\_05\_median} & 2.2367 & 2.2728 & 0.1562 & 147\\
\texttt{RLLR} & \texttt{duration\_s} & 2.4524 & 2.0326 & 0.2006 & 110\\
\texttt{RLRL} & \texttt{prosody\_f0\_median} & 270.2333 & 2.0442 & 0.1015 & 320\\
\texttt{RLRR} & \texttt{prosody\_f0\_std} & 46.8924 & 1.9399 & 0.1148 & 296\\
\texttt{RRLL} & \texttt{mfcc\_03\_max} & 39.2929 & 1.4878 & 0.1099 & 272\\
\texttt{RRLR} & \texttt{mfcc\_d2\_10\_std} & 0.6478 & 0.7657 & 0.0986 & 260\\
\texttt{RRRL} & \texttt{mfcc\_d2\_00\_std} & 3.2140 & 2.1218 & 0.2554 & 106\\
\texttt{RRRR} & \texttt{mfcc\_d2\_12\_std} & 0.5605 & 1.6136 & 0.0943 & 329\\
\bottomrule
\end{tabular}
\end{table}

\begin{table}[H]
\caption{Entropy tree fitted on \textbf{MP4/AAC @ 64 kbit/s}, all internal nodes to depth 4.}
\label{tab:app-splits-mp4-aac64}
\centering\scriptsize\setlength{\tabcolsep}{1.5pt}
\begin{tabular}{@{}llrrrr@{}}
\toprule
Path & Feature & Threshold & $H$ & IG & $n$\\
\midrule
\textit{root} & \texttt{prosody\_intensity\_std} & 8.3241 & 2.5850 & 0.3064 & 4905\\
\texttt{L} & \texttt{mfcc\_00\_iqr} & 60.5853 & 2.3887 & 0.1408 & 3075\\
\texttt{R} & \texttt{mfcc\_d1\_00\_std} & 7.1056 & 2.0879 & 0.1555 & 1830\\
\texttt{LL} & \texttt{spectral\_band\_500\_1000\_mean} & 0.0797 & 1.9819 & 0.0875 & 1181\\
\texttt{LR} & \texttt{prosody\_f0\_mean} & 219.0372 & 2.4096 & 0.0957 & 1894\\
\texttt{RL} & \texttt{prosody\_f0\_median} & 161.0991 & 2.2673 & 0.1047 & 874\\
\texttt{RR} & \texttt{prosody\_intensity\_std} & 11.2496 & 1.6245 & 0.0884 & 956\\
\texttt{LLL} & \texttt{prosody\_f0\_mean} & 169.4067 & 1.4798 & 0.0877 & 395\\
\texttt{LLR} & \texttt{prosody\_f3\_bw\_median} & 712.6446 & 2.0975 & 0.0757 & 786\\
\texttt{LRL} & \texttt{spectral\_band\_0\_500\_mean} & 0.8284 & 2.3465 & 0.0696 & 1662\\
\texttt{LRR} & \texttt{mfcc\_02\_std} & 14.7880 & 2.0711 & 0.1528 & 232\\
\texttt{RLL} & \texttt{spectral\_band\_1000\_2000\_std} & 0.0765 & 2.3608 & 0.1810 & 226\\
\texttt{RLR} & \texttt{spectral\_band\_0\_500\_mean} & 0.7133 & 2.0917 & 0.0718 & 648\\
\texttt{RRL} & \texttt{spectral\_contrast\_2\_mean} & 12.2765 & 1.8882 & 0.1074 & 461\\
\texttt{RRR} & \texttt{mfcc\_03\_max} & 38.0092 & 1.2059 & 0.0851 & 495\\
\texttt{LLLL} & \texttt{mfcc\_05\_iqr} & 7.6452 & 1.3054 & 0.1059 & 241\\
\texttt{LLLR} & \texttt{prosody\_f2\_std} & 326.8660 & 1.5276 & 0.1362 & 154\\
\texttt{LLRL} & \texttt{prosody\_f0\_mean} & 199.2747 & 2.1863 & 0.1157 & 498\\
\texttt{LLRR} & \texttt{spectral\_contrast\_1\_mean} & 7.7606 & 1.7295 & 0.0923 & 288\\
\texttt{LRLL} & \texttt{duration\_s} & 2.8480 & 2.2806 & 0.0690 & 1204\\
\texttt{LRLR} & \texttt{prosody\_cpps} & 5.3755 & 2.2671 & 0.0731 & 458\\
\texttt{LRRL} & \texttt{mfcc\_13\_kurt} & 0.3511 & 1.5158 & 0.2506 & 66\\
\texttt{LRRR} & \texttt{spectral\_band\_5000\_6000\_std} & 0.0311 & 2.0779 & 0.1904 & 166\\
\texttt{RLLL} & \texttt{mfcc\_00\_iqr} & 112.3155 & 2.2882 & 0.1462 & 170\\
\texttt{RLLR} & \texttt{spectral\_band\_500\_1000\_mean} & 0.1727 & 1.8385 & 0.2681 & 56\\
\texttt{RLRL} & \texttt{prosody\_f0\_median} & 270.2040 & 2.0207 & 0.1057 & 296\\
\texttt{RLRR} & \texttt{prosody\_f0\_slope\_abs\_mean} & 704.7569 & 2.0192 & 0.1053 & 352\\
\texttt{RRLL} & \texttt{mfcc\_05\_median} & 3.1128 & 2.1051 & 0.1682 & 126\\
\texttt{RRLR} & \texttt{spectral\_band\_0\_500\_std} & 0.2993 & 1.6568 & 0.1037 & 335\\
\texttt{RRRL} & \texttt{prosody\_f4\_median} & 3950.6527 & 1.5464 & 0.1841 & 122\\
\texttt{RRRR} & \texttt{spectral\_band\_3000\_4000\_mean} & 0.0143 & 0.9815 & 0.0752 & 373\\
\bottomrule
\end{tabular}
\end{table}

\begin{table}[H]
\caption{Entropy tree fitted on \textbf{Roundtrip WAV (AAC re-wrapped)}, all internal nodes to depth 4.}
\label{tab:app-splits-roundtrip-wav}
\centering\scriptsize\setlength{\tabcolsep}{1.5pt}
\begin{tabular}{@{}llrrrr@{}}
\toprule
Path & Feature & Threshold & $H$ & IG & $n$\\
\midrule
\textit{root} & \texttt{prosody\_intensity\_std} & 8.3241 & 2.5850 & 0.3064 & 4905\\
\texttt{L} & \texttt{mfcc\_00\_iqr} & 60.5802 & 2.3887 & 0.1408 & 3075\\
\texttt{R} & \texttt{mfcc\_d1\_00\_std} & 7.1058 & 2.0879 & 0.1555 & 1830\\
\texttt{LL} & \texttt{spectral\_band\_500\_1000\_mean} & 0.0782 & 1.9819 & 0.0882 & 1181\\
\texttt{LR} & \texttt{prosody\_f0\_mean} & 219.0372 & 2.4096 & 0.0957 & 1894\\
\texttt{RL} & \texttt{prosody\_f0\_median} & 161.0993 & 2.2673 & 0.1047 & 874\\
\texttt{RR} & \texttt{prosody\_intensity\_std} & 11.2497 & 1.6245 & 0.0884 & 956\\
\texttt{LLL} & \texttt{prosody\_f0\_mean} & 170.6874 & 1.4582 & 0.0902 & 381\\
\texttt{LLR} & \texttt{prosody\_f3\_bw\_median} & 712.5267 & 2.0961 & 0.0727 & 800\\
\texttt{LRL} & \texttt{spectral\_band\_0\_500\_mean} & 0.8280 & 2.3465 & 0.0701 & 1662\\
\texttt{LRR} & \texttt{mfcc\_02\_std} & 14.7881 & 2.0711 & 0.1528 & 232\\
\texttt{RLL} & \texttt{spectral\_band\_1000\_2000\_std} & 0.0765 & 2.3608 & 0.1810 & 226\\
\texttt{RLR} & \texttt{spectral\_band\_0\_500\_mean} & 0.7136 & 2.0917 & 0.0728 & 648\\
\texttt{RRL} & \texttt{spectral\_contrast\_2\_mean} & 12.2826 & 1.8882 & 0.1051 & 461\\
\texttt{RRR} & \texttt{mfcc\_03\_max} & 42.9480 & 1.2059 & 0.0846 & 495\\
\texttt{LLLL} & \texttt{mfcc\_02\_min} & -4.2720 & 1.3393 & 0.0935 & 239\\
\texttt{LLLR} & \texttt{spectral\_flatness\_median} & 0.0075 & 1.4163 & 0.1184 & 142\\
\texttt{LLRL} & \texttt{prosody\_f0\_mean} & 199.2748 & 2.1897 & 0.1143 & 505\\
\texttt{LLRR} & \texttt{spectral\_contrast\_1\_mean} & 7.7596 & 1.7311 & 0.0940 & 295\\
\texttt{LRLL} & \texttt{duration\_s} & 2.8480 & 2.2798 & 0.0680 & 1203\\
\texttt{LRLR} & \texttt{prosody\_f1\_bw\_median} & 281.4372 & 2.2671 & 0.0744 & 459\\
\texttt{LRRL} & \texttt{mfcc\_13\_kurt} & 0.3489 & 1.5158 & 0.2506 & 66\\
\texttt{LRRR} & \texttt{spectral\_band\_5000\_6000\_std} & 0.0311 & 2.0779 & 0.1904 & 166\\
\texttt{RLLL} & \texttt{mfcc\_00\_iqr} & 112.3050 & 2.2882 & 0.1462 & 170\\
\texttt{RLLR} & \texttt{spectral\_band\_500\_1000\_mean} & 0.1727 & 1.8385 & 0.2681 & 56\\
\texttt{RLRL} & \texttt{prosody\_f0\_median} & 270.2047 & 2.0178 & 0.1036 & 299\\
\texttt{RLRR} & \texttt{prosody\_f0\_slope\_abs\_mean} & 704.7769 & 2.0200 & 0.1026 & 349\\
\texttt{RRLL} & \texttt{spectral\_hf\_ratio\_4000\_std} & 0.0443 & 2.1003 & 0.1608 & 128\\
\texttt{RRLR} & \texttt{spectral\_band\_0\_500\_std} & 0.3006 & 1.6592 & 0.1006 & 333\\
\texttt{RRRL} & \texttt{spectral\_band\_0\_500\_mean} & 0.6868 & 1.5345 & 0.1361 & 210\\
\texttt{RRRR} & \texttt{spectral\_band\_3000\_4000\_mean} & 0.0144 & 0.8168 & 0.1042 & 285\\
\bottomrule
\end{tabular}
\end{table}



\end{document}
